\section*{Exercises}
\begin{exercise}\label{ex:2.1}State the converse and contrapositive of ``If $x=3$, then $x^2=9$,'' for real $x$. Decide which are true.\end{exercise}
\begin{exercise}\label{ex:2.2}Negate: ``For every integer $n$, there is an integer $m$ with $m>n$ and $m$ even.'' Preserve the domain and negate the conjunction correctly.\end{exercise}
\begin{exercise}\label{ex:2.3}Decide whether $\forall x\in\RR\ \exists y\in\RR:\ x+y=0$ and $\exists y\in\RR\ \forall x\in\RR:\ x+y=0$ are true. Supply a witness or a refutation.\end{exercise}
\begin{exercise}\label{ex:2.4}Refute: ``If a real number satisfies $x^2>4$, then $x>2$.'' Repair the conclusion without restricting the domain.\end{exercise}
\begin{exercise}\label{ex:2.5}Two students each accept at least one task, but cannot receive distinct acceptable tasks. Give an example. What stronger condition would detect its failure?\end{exercise}
\begin{exercise}\label{ex:2.6}For each positive $\eps$, choose a positive integer $N$ such that every integer $n\ge N$ satisfies $1/n<\eps$. Explain both why $N$ may depend on $\eps$ and why it must not depend on the later input $n$.\end{exercise}
